\section{Discussion}
\label{sec:discussion}

The analysis in Section~\ref{sec:convergence} is intentionally learning-centered and relies on assumptions that should be read as modeling constraints rather than deployment guarantees. The prototype distributions use diagonal Gaussian class-wise prototype statistics, which keeps communication and memory lightweight but cannot represent all cross-feature dependencies. The bounded-drift and bounded-staleness assumptions in Assumptions~\ref{asm:drift} and~\ref{asm:staleness} are also restrictive: they make it possible to express the stationary-neighborhood radius in terms of prototype-induced distributional drift and stale prototype updates, but they do not cover arbitrary event failures or unbounded distribution shift. Similarly, the reliability-aware prototype aggregation rule in \eqref{eq:reliability_weight} is restricted to an analyzable monotone form; more expressive learned reliability functions would require additional structure before they could be inserted into Lemma~\ref{lem:reliability}.

The serverless framing should also be interpreted carefully. Serverless-style execution motivates stateless learning episodes, event loss, function timeout, cold-start delay, and partial edge execution, but the contribution of this manuscript is the \ehsfp learning algorithm and its convergence behavior. It is not a resource-allocation, scheduling, placement, or cloud-cost optimization method. Those system-level questions are orthogonal to the learning problem studied here: how edge and cloud models behave when they optimize over incomplete, stale, and memory-replayed class-wise prototype statistics.

The empirical evidence in Section~\ref{sec:experiments} has three limitations. First, only the matched method comparisons use five seeds; the robustness, ablation, and scalability sweeps are single-seed, so small differences between their cells are not conclusive. Second, the journal-protocol \hsfp backbone performs well below its conference-reference accuracy, and both prototype-based methods trail the strongest weight-sharing baselines. The gains of \ehsfp are therefore established relative to \hsfp, not relative to the state of the art. Third, \ehsfp does not improve segmentation, and ImageNet-1K was not rerun under the 60-round journal protocol. Future work will consider richer covariance modeling beyond diagonal Gaussian prototypes, learned reliability functions with analyzable monotonicity or calibration constraints, and deployment-level optimization for scheduling, scaling, and resource placement as separate system contributions. The matched 200-round evaluation shows consistent endpoint gains over H-SFP across the reported classification, segmentation, scalability, and temperature settings.
